% ===================================================================
%  도표 제 15 호 — 시장. 먹을거리.
%
%  Traduction, faite en 2026, en coreen d'aujourd'hui, sur l'ido de
%  texto/io. Regles de la colonne : voir 10-dopyo-01.tex. Dix-huit
%  alineas, cent cinq renvois — le plus grand nombre du livret.
%
%  ET DIX-HUIT PAIRES SEULEMENT. Le tableau 13, plus court de douze
%  renvois, en coutait trente-trois. Le marche est la forme extreme
%  de ce que le tableau 14 avait montre : il ne fait presque que des
%  listes. La colonne telougoue avait sept inversions a faire sur ces
%  memes dix-huit paires — la liste est la meme pour toutes les
%  colonnes, puisque parigi.py ne lit que l'ido. Meme conclusion,
%  meme place dans le classement.
%
%  LE COREEN EST CHEZ LUI DANS UN MARCHE, ET LA LISTE LE MONTRE :
%  숫돌 la pierre a aiguiser, 광주리 le panier, 됫박 la mesure,
%  절구 le mortier, 옹기 la poterie, 젓갈 le poisson sale, 미역
%  l'algue, 배추 le chou, 무 le navet, 마늘 l'ail, 파 le poireau,
%  참깨 le sesame. Les poissons et les coquillages ont tous leur nom,
%  et le coreen distingue ce que le francais confond : 대구 la morue,
%  청어 le hareng, 홍합 la moule, 굴 l'huitre, 게 le crabe, 새우 la
%  crevette, 가재 l'ecrevisse.
%
%  LA COUPURE EST EUROPEENNE, comme au tableau 15 telougou et sur
%  exactement les memes objets : 아티초크 l'artichaut, 카망베르,
%  그뤼예르 les fromages, 헤이즐넛 les noisettes, 올리브 les olives,
%  꿩 le faisan — celui-la vit en Coree et garde son nom coreen. Ni
%  le fromage ni l'artichaut n'existent dans une cuisine coreenne de
%  1926 ; ils gardent le leur.
%
%  ET UN MOT JUSTE DE 1926 QU'ON N'ECRIT PAS EN 2026, QUATRIEME FOIS.
%  백정 traduit exactement « boucher » et c'est, en coreen
%  d'aujourd'hui, le nom d'une caste 천민 et une injure. La colonne
%  ecrit 고기 파는 사람 et 정육점, la fonction et le lieu — comme le
%  telougou ecrivait మాంసం అమ్మేవాడు pour eviter కసాయి. Deux langues
%  sans parente, deux mots differents, la meme raison et la meme
%  solution.
% ===================================================================

%%K t15-apar-1 apar
\VUsaut{4.0mm}
\VUtitre{15.0pt}{60}{도표 제 15 호}
\VUsaut{3.0mm}

%%K t15-tit-1 sub
\VUpk{11.4pt}{40}{시장. --- 먹을거리.}
\VUsaut{3.0mm}

%%K t15-01-1 p
1. --- 우리 끼니에 필요한 것을 대어 주는 장사꾼 대부분은 \VUgras{지붕 있는 시장}의
\VUgras{건물} \textsuperscript{(1)} 아래나 이웃 거리에 모인다.

%%K t15-02-1 p
2. --- \VUgras{돼지고기 파는 사람} \textsuperscript{(2)}은 — 그는 \VUgras{햄}
\textsuperscript{(3)}, \VUgras{순대} \textsuperscript{(4)}, \VUgras{소시지}
\textsuperscript{(5)}, \VUgras{창자 소시지} \textsuperscript{(6)}, \VUgras{비계}
\textsuperscript{(7)}를 판다 — \VUgras{버터와 치즈 파는 여자} \textsuperscript{(8)} 곁에
선다. 그 여자의 좌판에는 붉은 껍질의 \VUgras{네덜란드 치즈} \textsuperscript{(9)},
\VUgras{카망베르} \textsuperscript{(10)}, \VUgras{그뤼예르 덩이} \textsuperscript{(11)},
싱싱한 \VUgras{버터 덩이} \textsuperscript{(12)}가 그득하다.

%%K t15-03-1 p
3. --- 그 여자 앞에서 \VUgras{남새 파는 여자} \textsuperscript{(13)}가 손님들에게 고운
\VUgras{토마토} \textsuperscript{(14)}, \VUgras{꽃양배추} \textsuperscript{(15)},
\VUgras{상추} \textsuperscript{(16)}, \VUgras{양배추} \textsuperscript{(17)},
\VUgras{호박} \textsuperscript{(19)}, \VUgras{참외} \textsuperscript{(18)}, 심지어
\VUgras{버섯} \textsuperscript{(20)}까지 내보인다. 그러나 \VUgras{맥주 나르는 사람}이 제
\VUgras{배달 수레} \textsuperscript{(21)}를 — \VUgras{맥주 통} \textsuperscript{(22)}을 실은
수레다 — 바로 그 좌판 앞에 세워서, 그 여자의 손님들이 다가오지 못한다. 밭일하는 여자
하나가 그 여자에게 \VUgras{아티초크} \textsuperscript{(23)} 광주리를 가져다준다.

%%K t15-tit-2 sub
\VUsaut{4.0mm}
\VUpk{10.2pt}{40}{팔기와 사기.}
\VUsaut{3.0mm}

%%K t15-04-1 p
4. --- 시장에는 활기가 넘친다. \VUgras{경매} \textsuperscript{(24)} 시간이 왔기 때문이다.
사러 오는 \VUgras{안주인} \textsuperscript{(26)}들은 \VUgras{내장 파는 여자}
\textsuperscript{(25)}의 가게 앞을 지나며 멈춘다 — \VUgras{양}이나 \VUgras{송아지 머리}를
사려는 것이다.

%%K t15-05-1 p
5. --- 살진 \VUgras{요리사} \textsuperscript{(27)} 하나가 아주 좋은 \VUgras{다랑어}를
흥정한다 — \VUgras{생선 파는 여자} \textsuperscript{(28)}에게서. 그 여자는 제멋대로 값을
올린다. 그 여자에게는 \VUgras{뱀장어}, \VUgras{가자미}, \VUgras{송어}, \VUgras{잉어}가
많이 들어왔다. 그 여자의 \VUgras{대구}와 \VUgras{홍합}도 아주 싱싱하다.

%%K t15-tit-3 sub
\VUsaut{4.0mm}
\VUpk{10.2pt}{40}{싼 것과 비싼 것.}
\VUsaut{3.0mm}

%%K t15-06-1 p
6. --- 그 곁에 자리 잡은 \VUgras{닭 파는 여자} \textsuperscript{(29)}는 아주 정직하다.
\VUgras{칠면조}, \VUgras{거위}, \VUgras{오리}, 또는 수수한 \VUgras{영계}를 팔 때 그
여자는 바른 값만 부른다.

%%K t15-07-1 p
7. --- 광장 모퉁이 이 층에는 얼마 전부터 \VUgras{옷감 상인} \textsuperscript{(30)}이
자리를 잡았다. 그에게서 사는 여자들은 늘 고운 \VUgras{상보}, \VUgras{손수건},
\VUgras{앞치마}, \VUgras{행주}를 넉넉히 찾을 수 있다. 같은 층에 \VUgras{칼 만드는 이}
\textsuperscript{(31)}가 제 일터를 차렸다. 그는 시장 여자들의 \VUgras{칼}을 갈아 주고,
밭일하는 이들에게는 가장 단단한 \VUgras{강철}로 \VUgras{낫}을 만들어 준다.

%%K t15-tit-4 sub
\VUsaut{4.0mm}
\VUpk{10.2pt}{40}{식료품.}
\VUsaut{3.0mm}

%%K t15-08-1 p
8. --- 그 일터 아래에 얼마 전부터 큰 식료품 가게가 열렸다. 그것은 이제 옛날의 수수하고
대수롭지 않은 \VUgras{식료품 가게} \textsuperscript{(32)}가 아니다 — 그 가게는 흔한 것만
팔았다. \VUgras{차} \textsuperscript{(33)}, \VUgras{초콜릿} \textsuperscript{(34)},
\VUgras{덩이 설탕} \textsuperscript{(37)}, \VUgras{비누} \textsuperscript{(38)}였다.
여기에서는 무엇이든 판다. \VUgras{바삭한 과자}, \VUgras{통조림} \textsuperscript{(35)},
\VUgras{술} \textsuperscript{(36)}, \VUgras{럼}, \VUgras{코냑}, \VUgras{키르슈},
\VUgras{아니스 술}, \VUgras{퀴라소}다. 여기에서 사냥한 고기도 찾을 수 있다.
\VUgras{산토끼} \textsuperscript{(39)}, \VUgras{개똥지빠귀} \textsuperscript{(40)},
\VUgras{자고} \textsuperscript{(41)}, \VUgras{메추라기} \textsuperscript{(42)},
\VUgras{들오리} \textsuperscript{(43)}, \VUgras{꿩} \textsuperscript{(44)}이다 — 그리고
\VUgras{바나나} \textsuperscript{(46)}나 \VUgras{파인애플} 같은 남쪽 과일도 그만큼 쉽게
찾을 수 있다.

%%K t15-09-1 p
9. --- 아주 친절한 식료품 가게 \VUgras{사환} \textsuperscript{(45)}은 원하는 것을 내줄
채비가 되어 있다. \VUgras{잼} \textsuperscript{(47)} 넉 되나 닷 되, \VUgras{기름}
\textsuperscript{(48)}, \VUgras{오이 절임} \textsuperscript{(49)}, 소금에 절인
\VUgras{정어리} \textsuperscript{(50)}다.

%%K t15-09-2 p
그것이 \VUgras{손님} \textsuperscript{(51)}들에게 아주 편하다. 가장 흔한 것 곁에서 가장
드문 첫물과 가장 맛있는 과일을 함께 찾을 수 있기 때문이다. 물 많은 \VUgras{배}
\textsuperscript{(52)}, \VUgras{자두} \textsuperscript{(53)}, \VUgras{포도}
\textsuperscript{(54)}, \VUgras{복숭아} \textsuperscript{(55)}, \VUgras{아몬드}
\textsuperscript{(56)}, \VUgras{호두} \textsuperscript{(57)}다.

%%K t15-tit-5 sub
\VUsaut{4.0mm}
\VUpk{10.2pt}{40}{손님들.}
\VUsaut{3.0mm}

%%K t15-10-1 p
10. --- \VUgras{쌀} \textsuperscript{(58)}과 \VUgras{렌즈콩} \textsuperscript{(59)}은
훈제한 \VUgras{청어} \textsuperscript{(60)} 궤 곁에 있다. \VUgras{감자} \textsuperscript{(61)}
와 \VUgras{강낭콩} \textsuperscript{(63)}은 \VUgras{사과} \textsuperscript{(62)},
\VUgras{무화과} \textsuperscript{(64)}, \VUgras{말린 자두} \textsuperscript{(65)},
\VUgras{헤이즐넛} \textsuperscript{(66)} 곁에 있다. 그 가게에서는 싱싱한 \VUgras{달걀}
\textsuperscript{(67)}과 \VUgras{올리브} \textsuperscript{(68)}도 찾을 수 있다. 그래서
\VUgras{광주리} \textsuperscript{(70)}와 \VUgras{장 그물망} \textsuperscript{(73)}이 아주
빠르게 채워진다.

%%K t15-11-1 p
11. --- 그 훌륭한 가게의 \VUgras{커피} \textsuperscript{(72)}는 \VUgras{하녀}
\textsuperscript{(69)}들과 \VUgras{찬모}들 사이에서 마땅한 이름을 얻었다. 늙은
\VUgras{식료품 상인} \textsuperscript{(71)} 자신이 그것을 아주 정성껏 볶기 때문이다.

%%K t15-tit-6 sub
\VUsaut{4.0mm}
\VUpk{10.2pt}{40}{볼거리.}
\VUsaut{3.0mm}

%%K t15-12-1 p
12. --- 모두가 물건을 사러 시장에 오는 것은 아니다. 시장 건물로 가는 좁은 길의 그림 같은
오감 속에서 아주 여러 사람이 보인다. 상냥한 \VUgras{정육점 사환} \textsuperscript{(75)}
곁에 — 그는 제 앞에 \VUgras{손수레} \textsuperscript{(74)}를 밀고 간다 — 휴가 나온 병사
둘이 있다. 제 번쩍이는 군복을 아주 자랑스레 보이는 사람들이다. 푸른 \VUgras{돌망}과
\VUgras{깃털 군모}를 쓴 \VUgras{경기병} \textsuperscript{(76)}, 그리고 부푼 소매에
\VUgras{페즈}를 쓴 \VUgras{주아브 하사}다.

%%K t15-13-1 p
13. --- \VUgras{빵 굽는 사람} \textsuperscript{(80)}은 — 그는 \VUgras{빵집} \textsuperscript{(78)}
문지방에 서 있다 — 방금 제 \VUgras{빵} \textsuperscript{(79)} 반죽을 치대고, 화덕에서
\VUgras{크루아상} \textsuperscript{(81)}과 \VUgras{작은 빵} \textsuperscript{(82)}과
\VUgras{둥근 빵} \textsuperscript{(83)}을 꺼냈다. 그것들이 지금 진열창에 있다.

%%K t15-14-1 p
14. --- 빵집 위에는 솜씨 좋은 \VUgras{속옷 짓는 여자} \textsuperscript{(84)}가 산다. 그
여자는 \VUgras{수놓는 여자} \textsuperscript{(85)}를 여럿 두고 있다. 그 곁에는
\VUgras{빨래하고 다림질하는 여자} \textsuperscript{(86)}가 산다. 그 여자는 \VUgras{빨래}
\textsuperscript{(88)}를 빨아 말린 뒤 풀을 먹이고 \VUgras{다리미} \textsuperscript{(87)}로
편다.

%%K t15-tit-7 sub
\VUsaut{4.0mm}
\VUpk{10.2pt}{40}{고기, 생선, 남새.}
\VUsaut{3.0mm}

%%K t15-15-1 p
15. --- 아래층에 있는 \VUgras{정육점} \textsuperscript{(89)}에서는 온갖 \VUgras{고기}를
판다 — \VUgras{양고기} \textsuperscript{(90)}, \VUgras{쇠고기} \textsuperscript{(94)},
\VUgras{송아지 고기} \textsuperscript{(92)}다 —, \VUgras{양 넓적다리},
\VUgras{비프스테이크}, \VUgras{구이}, \VUgras{커틀릿}의 꼴로. \VUgras{고기 써는 사람}
\textsuperscript{(93)}이 그 고기를 모두 토막 내고, \VUgras{골 든 뼈} \textsuperscript{(96)}
를 따로 놓는다. 정육점 문 앞에는 \VUgras{굴 파는 여자} \textsuperscript{(97)}가 있어
\VUgras{굴} \textsuperscript{(98)}을 판다. 원하면 그 여자가 그것을 까 준다.

%%K t15-16-1 p
16. --- 그 장사꾼들은 모두 세나 자릿세를 낸다. 그래서 \VUgras{순경} \textsuperscript{(100)}
은 딱한 \VUgras{떠돌이 남새 장수} \textsuperscript{(99)}들을 사정없이 쫓는다. 그 여자들이
제 작은 수레에 \VUgras{당근}, \VUgras{무}, \VUgras{순무}, \VUgras{파},
\VUgras{아스파라거스 묶음}, \VUgras{풋완두}, \VUgras{시금치}, \VUgras{수영},
\VUgras{물냉이}, \VUgras{파슬리}를 싣고 다니며 그들과 겨루기 때문이다.

%%K t15-tit-8 sub
\VUsaut{4.0mm}
\VUpk{10.2pt}{40}{순경을 조심하라!}
\VUsaut{3.0mm}

%%K t15-17-1 p
17. --- \VUgras{마늘} \textsuperscript{(101)}과 \VUgras{양파}를 파는 아이도 잘 안다 —
저 또한 시장 가까이에서 제 물건을 팔면 안 된다는 것을. 그는 달아난다. \VUgras{바닷가재}와
\VUgras{큰 새우}를 \VUgras{파는 사람} \textsuperscript{(102)}의 \VUgras{게}와
\VUgras{가재}와 \VUgras{새우} 광주리를 엎을 위험을 무릅쓰고 달리면서.

%%K t15-18-1 p
18. --- 모두가 움직이며 크게 떠든다. 그래도 돌아다니며 파는 \VUgras{유리 끼우는 사람}
\textsuperscript{(103)}의 목소리와 \VUgras{숯 파는 사람} \textsuperscript{(104)}의 외침은
또렷이 들린다 — 그 숯장수는 \VUgras{숯 한 자루} \textsuperscript{(105)}를 배달하려고 제
수레를 잠시 두고 왔다.

\VUsaut{6.0mm}
\VUfilet{22mm}
